\chapter{Your first complete database}
\label{learn:tutorial}
\index{tutorial}
\learning{Create a fresh database, load a relational schema, insert data, retrieve useful answers, make reversible changes and verify the results. Every file in this chapter is supplied with the manual.}

\section{Prepare two terminals and one fresh database}
Complete the installation chapter first. Work from the repository root. Use
one terminal for the server and another for the client. In the first terminal
create a new teaching database:
\begin{lstlisting}[language=bash]
./build/bin-linux/minisqld --init ./data shop 4096
\end{lstlisting}
On Windows the equivalent command is:
\begin{lstlisting}[language=PowerShell]
.\build\bin\minisqld.exe --init .\data shop 4096
\end{lstlisting}
The output includes a directory beginning with \code{db\_}, a human-readable
database name and a page size. Copy the actual directory into the server
command. The following placeholder is deliberately not a real directory:
\begin{lstlisting}[language=bash]
./build/bin-linux/minisqld --serve ./data/db_<uuid> 7432 16
\end{lstlisting}
The server remains running. A listener-ready message means the server has
opened the database and reached its listener setup. Check the complete
startup output if the next command cannot connect.

In the second terminal, test the connection:
\begin{lstlisting}[language=bash]
./build/bin-linux/minisql --ping 7432
./build/bin-linux/minisql --shell 7432
\end{lstlisting}
\code{--ping} prints \code{PONG} on success. The shell displays a
\code{minisql>} prompt. The trusted listener is a local learning mode: its
sessions have administrative authority. Read the security chapter before
allowing users or applications to connect in a deployment.

\section{Understand the model before typing the schema}
The teaching database represents customers, products, orders and individual
order lines. A customer may have no orders. A product may have stock but no
orders. One order can contain several products. The schema retains these
distinctions instead of trying to fit all facts into one table.
\begin{figure}[htbp]
\centering
\begin{tikzpicture}[node distance=18mm and 20mm]
\node[strong,text width=55mm,align=left] (customers) {\textbf{customers}\\\rule{54mm}{.3pt}\\\textbf{customer\_id} : INTEGER (PK)\\name : VARCHAR(80)\\city : VARCHAR(60)\\active : BOOLEAN};
\node[strong,text width=55mm,align=left,right=of customers] (orders) {\textbf{orders}\\\rule{54mm}{.3pt}\\\textbf{order\_id} : INTEGER (PK)\\customer\_id : INTEGER (FK)\\order\_status : VARCHAR(20)};
\node[strong,text width=55mm,align=left,below=of customers] (products) {\textbf{products}\\\rule{54mm}{.3pt}\\\textbf{product\_id} : INTEGER (PK)\\name : VARCHAR(100)\\price : DECIMAL(10,2)\\stock : INTEGER};
\node[strong,text width=55mm,align=left,below=of orders] (items) {\textbf{order\_items}\\\rule{54mm}{.3pt}\\\textbf{order\_id, product\_id} (PK)\\Both also foreign keys\\quantity : INTEGER\\unit\_price : DECIMAL(10,2)};
\draw[arrow] (customers) -- node[smalllabel,above] {1 to many} (orders);
\draw[arrow] (orders) -- node[smalllabel,right] {1 to many} (items);
\draw[arrow] (products) -- node[smalllabel,above] {1 to many} (items);
\end{tikzpicture}
\caption{The teaching database: one customer can place many orders; order items connect orders to products. Keys are immutable in these examples.}
\label{fig:shop-schema}
\end{figure}


The examples use explicit stable identifiers so that expected results are
easy to compare. Identity and sequence features are taught later, with their
different allocation semantics. The identifier of an existing customer or
product is never changed during this tutorial.

\section{Create and populate the tables}
Leave the shell with \code{\textbackslash quit}, or open a third terminal,
and run the first script:
\begin{lstlisting}[language=bash]
./build/bin-linux/minisql --script 7432 \
  docs/manual/examples/01-create-shop.sql
\end{lstlisting}
On Windows:
\begin{lstlisting}[language=PowerShell]
.\build\bin\minisql.exe --script 7432 `
  .\docs\manual\examples\01-create-shop.sql
\end{lstlisting}
The file contains the following complete setup. It is intentionally designed
for a new empty database. Rerunning it on the same database should fail when
objects already exist; an unexpected pre-existing schema must not be silently
accepted as the tutorial schema.
\lstinputlisting[language=SQL,caption={The complete reproducible teaching database}]{examples/01-create-shop.sql}

The four \code{CREATE TABLE} statements are schema changes. They run before
the explicit transaction that loads data. The \code{BEGIN}/\code{COMMIT}
pair makes the initial row population one transaction. Parents are inserted
before children, so foreign-key checks can find the referenced rows.
\code{CHECK (stock >= 0)} prevents a negative quantity; \code{CHECK (quantity
> 0)} requires a positive order-line quantity. \code{NOT NULL} requires a
value rather than leaving it unknown.

The final count queries should report the following values:
\begin{table}[htbp]\centering
\caption{The teaching database after successful setup}
\begin{tabular}{l r l}\toprule
Table & Rows & What one row means\\\midrule
customers & 4 & One customer\\
products & 5 & One product offered for sale\\
orders & 4 & One order placed by a customer\\
order\_items & 6 & One product line in an order\\\bottomrule
\end{tabular}
\end{table}

\tip{What to do after a setup error}{Read the first error, not only the final
script status. A script stops at its failure point; earlier DDL may already
be committed. For the tutorial, create another fresh database instead of
guessing which objects to delete. The validation tool always uses a new
temporary directory.}

\section{Inspect rows and choose columns}
Reconnect with \code{--shell 7432} and run:
\begin{lstlisting}[language=SQL]
SELECT customer_id, name, city, active
FROM customers
ORDER BY customer_id;
\end{lstlisting}
\begin{table}[htbp]\centering
\caption{The four customers in identifier order}
\begin{tabular}{r l l l}\toprule
customer\_id & name & city & active\\\midrule
1 & Ada & London & TRUE\\
2 & Grace & New York & TRUE\\
3 & Linus & Helsinki & TRUE\\
4 & Margaret & NULL & FALSE\\\bottomrule
\end{tabular}
\end{table}
The selected column list controls the result's shape. It can omit columns,
reorder them or add expressions. A column alias such as
\code{COUNT(*) AS customer\_count} gives the result a readable heading.
Use explicit columns in application code so a later schema extension does
not accidentally change the expected result layout.

\section{Filter and sort}
\begin{lstlisting}[language=SQL]
SELECT product_id, name, stock
FROM products
WHERE stock < 15
ORDER BY stock, product_id;
\end{lstlisting}
The result contains Monitor Stand (104) with stock 8 and Mechanical Keyboard
(101) with stock 12, in that order. \code{WHERE} keeps rows for which its
condition is true. \code{ORDER BY stock, product\_id} sorts by stock, then
uses the identifier as a deterministic tie-breaker.

To find the two products with the most stock, reverse the first sort:
\begin{lstlisting}[language=SQL]
SELECT product_id, name, stock
FROM products
ORDER BY stock DESC, product_id
LIMIT 2;
\end{lstlisting}
The result is USB-C Cable (103, stock 100) followed by Notebook (105, stock
50). A bare \code{LIMIT 2} without \code{ORDER BY} would mean ``any first two
rows produced by the chosen plan'', not ``the two largest''.

\section{Work correctly with unknown values}
\begin{lstlisting}[language=SQL]
SELECT customer_id, name FROM customers
WHERE city IS NULL;

SELECT COUNT(city) AS known_cities,
       COUNT(*) AS customer_count
FROM customers;
\end{lstlisting}
Margaret is the customer with no recorded city. The counts are 3 known cities
and 4 customers. \code{COUNT(city)} counts non-null city values;
\code{COUNT(*)} counts rows. The test for an unknown value is
\code{IS NULL}, not \code{= NULL}. The language chapter develops the
three-valued logic behind this distinction.

\warning{NULL in the current text protocol}{The native console displays SQL
NULL as \code{NULL}. Protocol version 1 also uses that text for an actual
string containing the four characters \code{NULL}. The general Python/JDBC
result path consequently cannot distinguish them. Do not choose that literal
string as a special missing-data marker. This limitation is separate from
the engine's internal typed null representation.}

\section{Join related tables}
\begin{lstlisting}[language=SQL]
SELECT o.order_id, c.name, o.order_status
FROM orders o
JOIN customers c ON c.customer_id = o.customer_id
ORDER BY o.order_id;
\end{lstlisting}
\code{o} and \code{c} are table aliases. Qualified names such as
\code{o.customer\_id} make each column's origin clear. The join condition
matches an order's customer key to a customer row.

\begin{table}[htbp]\centering
\caption{Joining orders to their customers}
\begin{tabular}{r l l}\toprule
order\_id & name & order\_status\\\midrule
1001 & Ada & new\\
1002 & Grace & shipped\\
1003 & Ada & paid\\
1004 & Linus & new\\\bottomrule
\end{tabular}
\end{table}
Notice that Ada appears twice because she has two orders. A join produces
matching row combinations; it does not implicitly deduplicate customer
names. If a customer has no orders, an inner join omits that customer.

\section{Retain unmatched rows and count correctly}
\begin{lstlisting}[language=SQL]
SELECT c.name, COUNT(o.order_id) AS order_count
FROM customers c
LEFT JOIN orders o ON o.customer_id = c.customer_id
GROUP BY c.name
ORDER BY c.name;
\end{lstlisting}
The left join retains every customer. Margaret's unmatched right-hand
columns are null, so \code{COUNT(o.order\_id)} returns zero for her group.
The four results are Ada 2, Grace 1, Linus 1 and Margaret 0. Replacing
\code{COUNT(o.order\_id)} with \code{COUNT(*)} would count the retained
unmatched row as one. This is a common beginner mistake.

Names happen to be unique in this small fixture. In a real customer report,
group by the stable customer identifier as well as the display name, or two
different people with the same name could be merged into one group.

\section{Build a grouped business query}
\begin{lstlisting}[language=SQL]
SELECT c.name, SUM(i.quantity) AS units
FROM customers c
JOIN orders o ON o.customer_id = c.customer_id
JOIN order_items i ON i.order_id = o.order_id
GROUP BY c.name
HAVING SUM(i.quantity) >= 3
ORDER BY c.name;
\end{lstlisting}
This asks for customers whose orders contain at least three units in total.
The result is Ada 6 and Grace 3. \code{SUM(i.quantity)} adds integer
quantities within each customer's group. \code{HAVING} filters groups after
aggregation. \code{WHERE}, if present, would filter the input rows before
grouping. The total quantity across all six order lines is 10.
\begin{figure}[htbp]
\centering
\begin{tikzpicture}[node distance=7mm and 16mm]
\node[strong,text width=57mm] (question) {\textbf{Business question}\\Which customers ordered at least three units?};
\node[strong,text width=57mm,right=of question] (physical) {\textbf{Physical choices}\\Observe with EXPLAIN};
\node[box,text width=57mm,below=of question] (sources) {FROM / JOIN\\Find related rows};
\node[box,text width=57mm,below=of physical] (access) {Index or sequential scan\\Join order and join algorithm};
\node[box,text width=57mm,below=of sources] (filter) {WHERE, then GROUP BY\\Filter rows and form groups};
\node[box,text width=57mm,below=of access] (blocking) {Filter pushdown / aggregation\\Memory or temporary spill};
\node[box,text width=57mm,below=of filter] (result) {HAVING / SELECT / ORDER BY\\Keep groups, choose columns, sort};
\node[box,text width=57mm,below=of blocking] (delivery) {Projection / sort / Top-N\\Bounded result frames};
\foreach \a/\b in {question/sources,sources/filter,filter/result,physical/access,access/blocking,blocking/delivery}{\draw[arrow] (\a) -- (\b);}
\draw[arrow,dashed] (sources) -- (access);\draw[arrow,dashed] (filter) -- (blocking);\draw[arrow,dashed] (result) -- (delivery);
\end{tikzpicture}
\caption{A learning model for a grouped query. SQL describes a result; the optimizer selects the physical work. This is a conceptual ordering, not a promise that every operator executes in the written order.}
\label{fig:query-thinking}
\end{figure}


\section{Interpret the price columns without assuming a wire format}
\begin{lstlisting}[language=SQL]
SELECT product_id, price FROM products WHERE product_id = 101;
\end{lstlisting}
At this revision the console returns \code{101 | 7990}. The declared scale is
2, so the stored coefficient 7990 represents the decimal value 79.90.
Exact decimal storage and exact decimal input validation are present, but
the generic wire formatter emits the coefficient without its decimal point.
The SQL arithmetic path also has a separate rescaling defect. For example,
the tested expression adding decimal 3.30 and 1.20 produces wire text 45000,
not a correct decimal sum representation.

The tutorial therefore uses integer quantities for SQL aggregation.
For an application that needs money calculations at this revision, one
workable design is to store and calculate integer minor units explicitly
(for example, cents), enforce a consistent scale in the application, and
format the result there. Range-check multiplication and sums. If reading an
existing \code{DECIMAL(10,2)} column, the application must know that column's
scale and decode its coefficient deliberately; ordinary result metadata does
not carry enough type information to reconstruct every value automatically.
Do not compensate for the arithmetic bug by guessing an extra divisor in
arbitrary SQL expressions.

\section{Make reversible changes}
Run the supplied transaction script:
\begin{lstlisting}[language=bash]
./build/bin-linux/minisql --script 7432 \
  docs/manual/examples/03-transactions.sql
\end{lstlisting}
Its first experiment is:
\begin{lstlisting}[language=SQL]
BEGIN;
UPDATE products SET stock = stock - 1 WHERE product_id = 101;
SAVEPOINT before_second;
UPDATE products SET stock = stock - 2 WHERE product_id = 101;
ROLLBACK TO before_second;
RELEASE before_second;
ROLLBACK;
\end{lstlisting}
The complete file queries the intermediate values. Stock starts at 12,
changes to 11, then 9, returns to 11 after the savepoint rollback, and returns
to 12 after the full rollback. A savepoint gives a named position within the
transaction; it does not commit earlier work. \code{RELEASE} removes that
savepoint. The final \code{ROLLBACK} abandons the transaction.

\begin{figure}[htbp]\centering
\begin{tikzpicture}[node distance=11mm and 11mm]
\node[box,text width=32mm] (initial) {Initial state\\stock = 12};
\node[strong,text width=32mm,right=of initial] (first) {First update\\stock = 11};
\node[box,text width=32mm,right=of first] (second) {Second update\\stock = 9};
\node[strong,text width=32mm,below=of first] (save) {Named savepoint\\remembers 11};
\node[box,text width=32mm,below=of initial] (rolled) {Full rollback\\stock = 12};
\draw[arrow] (initial) -- node[smalllabel,above] {BEGIN, update} (first);
\draw[arrow] (first) -- node[smalllabel,above] {update again} (second);
\draw[arrow] (first) -- node[smalllabel,right] {SAVEPOINT} (save);
\draw[arrow] (second.south) |- node[smalllabel,pos=.75,below] {ROLLBACK TO} (save.east);
\draw[arrow] (save) -- node[smalllabel,above] {ROLLBACK} (rolled);
\end{tikzpicture}
\caption{The tutorial's reversible stock experiment. Rolling back to a savepoint leaves the transaction active; a full rollback ends the transaction and restores its initial state.}
\label{fig:savepoints}
\end{figure}


Its second experiment inserts customer 99, updates the city, and deletes the
row, returning the changed values with \code{RETURNING}. A final rollback
leaves the teaching dataset unchanged. The check for customer 99 returns
zero rows counted. Read the transaction chapter before depending on failure,
DDL or concurrency behavior in application code.

\section{See constraints reject mistakes}
Run each of the following separately, outside an explicit transaction. Each
is an intentional error and is therefore not placed in the successful setup
script:
\begin{lstlisting}[language=SQL]
INSERT INTO customers(customer_id, name) VALUES (1, 'Duplicate');
INSERT INTO orders(order_id, customer_id, order_status)
VALUES (9999, 9999, 'new');
UPDATE products SET stock = -1 WHERE product_id = 101;
\end{lstlisting}
The first conflicts with an existing primary key. The second references a
customer that does not exist. The third violates the stock check. Query the
original rows afterwards to confirm that the rejected operations did not
perform the intended invalid change. In application code, treat a constraint
failure as a result to handle; do not convert every database error into
``try again forever''.

\section{Move to CTEs and window functions}
\path{examples/04-query-workshop.sql} gives named intermediate queries,
left-join anti-matching, per-customer order numbering, conditional display
values and a recursive sequence from 1 to 5. These build on the same schema.
The CTE that totals units per order returns 3, 3, 3 and 1 for orders
1001--1004. The window query numbers Ada's two orders 1 and 2 while restarting
at 1 for each other customer.

The query that finds customers without orders uses:
\begin{lstlisting}[language=SQL]
SELECT c.customer_id, c.name
FROM customers c
LEFT JOIN orders o ON o.customer_id = c.customer_id
WHERE o.order_id IS NULL
ORDER BY c.customer_id;
\end{lstlisting}
The result is Margaret (4). This uses an implemented, verified query shape.
The correlated-subquery section describes the narrower support and the
runtime issue found during validation; do not assume every standard
anti-join spelling is interchangeable in this development revision.

\section{Stop cleanly and reproduce your results}
When you have finished, issue \code{SHUTDOWN;} from the trusted shell. The
listener acknowledges the request, drains workers and closes its database.
You can then run offline tools on the directory. Reopen the same directory
with \code{--serve} to resume; do not run \code{--init} as a way of opening
an existing database.

To execute the tutorial in an isolated temporary database and compare its
results automatically, run:
\begin{lstlisting}[language=bash]
python3 docs/manual/tools/validate_examples.py \
  --bin-dir build/bin-linux \
  --output build/manual-validation.json
\end{lstlisting}
On Windows use \code{python}, the Windows binary directory, and normal
PowerShell line continuation if splitting the command. The runner starts
its own server on a free loopback port and removes its temporary database
afterwards. It does not open your teaching or application database.

\tip{You have completed the first learning milestone}{You can now create
and connect to a database, define relationships, load data, ask sorted and
grouped questions, retain unmatched rows, observe constraint errors and
undo a change. The next chapters explain each SQL feature systematically,
including syntax and MiniSQL-specific limitations.}
\sourcecode{src/minisql/sql/parser.ml}{The SQL parser and reserved-word behavior; the tutorial uses order\_status because STATUS is a keyword.}
\sourcecode{src/minisql/client/formatter.ml}{Result text formatting, including decimal coefficients.}
\sourcecode{src/minisql/executor/dml.ml}{Constraint checks and row changes used in the tutorial.}
